\section{Introduction}

Recently, video generation models have witnessed a remarkable progress, driven in particular by the rapid evolution of diffusion models \citep{ho2020denoising, song2020denoising, peebles2023scalable, rombach2022high} and VLM caption models \citep{liu2023visual, wang2024qwen2, team2024gemini}. 
Existing video generation models \citep{hong2022cogvideo, wan2025wan, kong2024hunyuanvideo, lin2024open} have achieved promising results in terms of visual fidelity and temporal consistency. 
Nonetheless, even state-of-the-art generators still face significant challenges in synthesizing complex human motions. The resulting videos often suffer from severe anatomical artifacts and unnatural joint articulations, as exemplified in Figure~\ref{fig:teaser}(a).
This deficiency primarily arises from the inherent limitations of pixel-only-regression training objectives, which tend to prioritize visual fidelity over the underlying kinematic principles governing human articulation. 
Concretely, the pixel-level reconstruction loss commonly used in diffusion models is dominated by static appearance and background details, not temporal kinetics. This deficiency is particularly pronounced for human subjects, whose high degrees of freedom make even subtle kinematic errors glaringly unnatural.

\begin{figure}[t!]
    \centering
    \includegraphics[width=\linewidth]{figures/teaser_.pdf}
    \caption{Overview of EchoMotion capabilities: (a) Improving anatomical integrity in human-centric video synthesis and (b) enabling bidirectional control between video and motion. By processing visual and motion sequences within a unified dual-branch Diffusion Transformer, the model learns a joint distribution of human appearance and kinematics.}
    \label{fig:teaser}
    \vspace{-9mm}
\end{figure}


Prior research~\citep{chefer2025videojam, huang2024vbench} has
highlighted that human video synthesis problems persist even with basic motion types that are well-represented in training data. This finding echoes the insight that the lack of anatomical plausibility is not merely a matter of data scale, but rather points to inherent difficulties in modeling fine-grained kinematic dynamics. To address this issue, some existing works focus on conditioning generation on explicit structural guidance, employing techniques like 2D keypoints priors\citep{jiang2025vace,ma2024follow} or 3D poses priors\citep{zhou2024realisdance, buchheim2025controlling}. 
While this approach offers direct control, it introduces two significant limitations. First, it creates a dependency on control signals that are often unavailable in real-world applications. Second, even when 3D priors like human poses are used, they are typically projected into the 2D image plane to align with the video frames. This projection process inevitably discards crucial 3D geometric information, leading to a diminished understanding of the underlying body structure and potential motion inconsistencies.

Motivated by these findings, we propose EchoMotion, an innovative framework that natively models the joint distribution of video and human motion modality.
It features a dual-modality diffusion transformer, which adopts a dual-branch architecture to uniformly process tokens from different modalities. Unlike the previous MMDiT architecture \citep{esser2024scaling}, which solely focuses on denoising video inputs, EchoMotion takes a further step by explicitly denoising parametric motion jointly. 
Benefiting from this explicit motion modeling, our method can significantly reduce motion artifacts and facilitate the generation of complex human motion videos.

In contrast to common approaches that condition on motion videos \citep{zhou2024realisdance, hu2024animate}, our parameterized motion representation is not only more token-efficient but also preserves native 3D structural information, which is more conducive to the model learning kinematic patterns.
Specifically, visual and motion tokens are concatenated along the sequence dimension and processed by our proposed Motion-Video Synchronized RoPE (MVS-RoPE) mechanism. This mechanism applies token-wise position embeddings to both modalities, enforcing precise temporal correspondence while preventing spatial positional collisions between the video and motion streams.
To ensure efficient model convergence and fully exploit the multi-task capabilities of our model, we design a Motion-Video Two-Stage Training Strategy.
This strategy employs a two-stage training recipe: an initial motion-only training phase followed by a \reb{motion-video multi-task training} phase. In the second phase, we further construct three paradigms: joint generation, motion-to-video generation, and video-to-motion generation, to efficiently capture the cross-modal interactions. Furthermore, we introduce a new high-quality, human-centric video dataset, which we name \textit{HuMoVe}. This dataset comprises approximately 80,000 video clips, each featuring distinct and clear human motions. For every video, we provide a granular textual description that details (1) the subject's appearance and attire, (2) the background context, and (3) a precise description of the action being performed. To enable multi-modal joint training, we also extract and include the corresponding SMPL motion parameters for each video clip, offering rich structural guidance.


Our experiments demonstrate that by holistically modeling the joint distribution of video and motion, EchoMotion significantly enhances the synthesis of human videos in terms of temporal coherence and structural integrity. Comprehensive evaluations validate its superiority over state-of-the-art baselines, showcasing drastic reductions in motion artifacts and superior preservation of physical plausibility. Furthermore, this unified approach inherently enables versatile cross-modal controllable generation, a capability confirmed through extensive qualitative and quantitative assessments.